\documentclass[11pt]{amsart}
\usepackage{amsmath,amssymb,amsthm}
\usepackage{graphicx}
\usepackage{hyperref}

\newtheorem{theorem}{Theorem}
\newtheorem{corollary}{Corollary}
\newtheorem{lemma}{Lemma}
\theoremstyle{definition}
\newtheorem{definition}{Definition}

\setlength{\emergencystretch}{3em}

\title[Irreducible-surjection extension theorem]{Irreducible-surjection conditions for unique order-preserving isometry extensions}
\author{Agent lane 16}
\date{2026-06-29}

\begin{document}
\maketitle

\section*{Status}
This is a candidate full literal answer to the final Problem in Tanmoy Paul and
T. S. S. R. K. Rao, \emph{Order-preserving unique Hahn-Banach extensions},
arXiv:2504.03386~\cite{PaulRao}.  The problem asks for conditions under which
an order-preserving isometric embedding
\[
  \Phi:A(K)\longrightarrow C(K)
\]
has a unique extension to an into order-preserving isometry
\[
  T:C(K)\longrightarrow C(K).
\]
The theorem below gives such conditions in the unital setting used throughout
the source paper: the state map associated with \(\Phi\) is required to be an
irreducible continuous surjection.  This is not a classification of every
possible embedding, but it is a direct set of conditions answering the source
request as stated.

\section*{Source problem}
The source problem appears on PDF page 14, in the final displayed Problem of
\cite{PaulRao}.  A crop of the statement is included below.

\begin{center}
\includegraphics[width=\linewidth]{figures/open_problem_crop.png}
\end{center}

The final sentence asks:
\begin{quote}
In view of the operator versions from Rao, it would be interesting to find
conditions under which an order-preserving isometric embedding
\(\Phi:A(K)\to C(K)\) has unique extension to an into order-preserving isometry
from \(C(K)\) to \(C(K)\).
\end{quote}

\section*{Definitions}
Let \(K\) be a compact convex subset of a locally convex real topological
vector space.  Let \(A(K)\) be the real Banach space of continuous affine
real-valued functions on \(K\), and let \(C(K)\) be the real Banach lattice of
continuous real-valued functions on \(K\), both with the supremum norm.

\begin{definition}
A continuous surjection \(\alpha:K\to K\) is called irreducible if no proper
closed subset \(L\subsetneq K\) satisfies \(\alpha(L)=K\).
\end{definition}

Every homeomorphism is irreducible, and the identity map is irreducible.
There are also non-injective irreducible surjections in general topology, so
the condition is broader than invertibility.

\section*{Idea of the Proof}
The source problem is an into-isometry problem for \(C(K)\)-spaces, so the
right structural input is Holsztynski's into form of the Banach-Stone theorem.
It says that any linear isometry \(T:C(K)\to C(K)\) is forced to be a
composition operator on some closed subset \(L\subset K\) that still sees the
whole domain space: there is a continuous surjection \(\beta:L\to K\) such
that \(Tf(t)=f(\beta(t))\) for \(t\in L\).  If \(T\) is positive and unital, no
sign factor appears.

For an order-preserving unital embedding \(\Phi:A(K)\to C(K)\), evaluation at
each \(t\in K\) gives a state of \(A(K)\), hence a point \(\alpha(t)\in K\);
so \(\Phi(a)(t)=a(\alpha(t))\).  If \(T\) extends \(\Phi\), then on the
Holsztynski set \(L\) we have
\[
  a(\beta(t))=Ta(t)=\Phi(a)(t)=a(\alpha(t))
\]
for every affine function \(a\).  Since affine functions separate points,
\(\beta(t)=\alpha(t)\) on \(L\).  But \(\beta(L)=K\), so
\(\alpha(L)=K\); irreducibility forces \(L=K\).  Thus \(T\) is the composition
operator \(f\mapsto f\circ \alpha\) everywhere.

\section*{The theorem}
\begin{lemma}[State-map form of unital positive embeddings]
Let \(\Phi:A(K)\to C(K)\) be a unital positive linear map.  Then there is a
unique continuous map \(\alpha_\Phi:K\to K\) such that
\[
  \Phi(a)(t)=a(\alpha_\Phi(t))\qquad(a\in A(K),\ t\in K).
\]
If \(\Phi\) is an isometry, then the closed convex hull of
\(\alpha_\Phi(K)\) is \(K\).  If \(\alpha_\Phi\) is onto, then \(\Phi\) is an
isometry.
\end{lemma}

\begin{proof}
For each \(t\in K\), the functional \(a\mapsto \Phi(a)(t)\) is positive and
takes the constant \(1\) to \(1\).  It is therefore a state of the order-unit
space \(A(K)\).  As recalled in the source paper's preliminaries, the state
space of \(A(K)\) is affinely homeomorphic to \(K\), so there is a unique
\(\alpha_\Phi(t)\in K\) such that
\(\Phi(a)(t)=a(\alpha_\Phi(t))\) for all \(a\in A(K)\).
Continuity of \(\alpha_\Phi\) follows because the functions in \(A(K)\)
separate points and each \(t\mapsto \Phi(a)(t)\) is continuous.

For \(a\in A(K)\),
\[
  \|\Phi(a)\|=\sup_{t\in K}|a(\alpha_\Phi(t))|
  =\sup_{x\in \alpha_\Phi(K)}|a(x)|.
\]
Thus \(\Phi\) is isometric exactly when \(\alpha_\Phi(K)\) is norming for
\(A(K)\).  By the Hahn-Banach separation theorem for compact convex sets, this
is equivalent to the closed convex hull of \(\alpha_\Phi(K)\) being \(K\).
In particular, if \(\alpha_\Phi\) is onto then \(\Phi\) is an isometry.
\end{proof}

\begin{theorem}
Let \(\Phi:A(K)\to C(K)\) be a unital order-preserving linear map, and let
\(\alpha=\alpha_\Phi\) be the associated state map.  Assume that
\(\alpha:K\to K\) is an irreducible continuous surjection.  Then \(\Phi\) is
an order-preserving isometric embedding, and it has a unique positive linear
isometric extension \(T:C(K)\to C(K)\).  This extension is
\[
  T(f)=f\circ \alpha\qquad(f\in C(K)).
\]
\end{theorem}

\begin{proof}
Since \(\alpha\) is onto, the lemma gives that \(\Phi\) is an isometry.  The
operator
\[
  C_\alpha:C(K)\to C(K),\qquad C_\alpha(f)=f\circ\alpha,
\]
is positive, unital, and isometric, and it extends \(\Phi\).  Thus existence
is immediate.

Let \(T:C(K)\to C(K)\) be any positive linear isometry extending \(\Phi\).
Since \(1\in A(K)\) and \(\Phi(1)=1\), we have \(T1=1\).  By Holsztynski's
theorem for into isometries of \(C(K)\)-spaces~\cite{Holsztynski}, there are a
closed subset \(L\subset K\) and a continuous surjection \(\beta:L\to K\) such
that
\[
  Tf(t)=f(\beta(t))\qquad(f\in C(K),\ t\in L).
\]
The usual unimodular sign factor in the real Banach-Stone representation is
\(1\) here because \(T1=1\).

For \(a\in A(K)\) and \(t\in L\),
\[
  a(\beta(t))=Ta(t)=\Phi(a)(t)=a(\alpha(t)).
\]
Continuous affine functions separate points of \(K\), hence
\(\beta(t)=\alpha(t)\) for every \(t\in L\).  Since \(\beta(L)=K\), it follows
that \(\alpha(L)=K\).  Irreducibility of \(\alpha\) forces \(L=K\).  Therefore,
for every \(f\in C(K)\) and \(t\in K\),
\[
  Tf(t)=f(\beta(t))=f(\alpha(t)).
\]
So \(T=C_\alpha\), proving uniqueness.
\end{proof}

\begin{corollary}[Canonical inclusion]
For every compact convex \(K\), the canonical unital inclusion
\[
  A(K)\subseteq C(K)
\]
has exactly one positive linear isometric extension \(C(K)\to C(K)\), namely
the identity operator.
\end{corollary}

\begin{proof}
Here \(\alpha=\mathrm{id}_K\), which is an irreducible continuous surjection.
The theorem applies.
\end{proof}

\begin{corollary}[Homeomorphic composition embeddings]
For every compact convex \(K\) and every homeomorphism \(h:K\to K\), the
embedding \(\Phi(a)=a\circ h\) from \(A(K)\) into \(C(K)\) has the unique
positive linear isometric extension \(T(f)=f\circ h\).  No metrizability or
point-peak hypothesis is needed.
\end{corollary}

\section*{Verification and novelty notes}
The proof is analytic and uses no numerical computation.  The decisive
external ingredient is Holsztynski's into-isometry theorem for spaces of
continuous functions; the rest is the standard state-space representation of
\(A(K)\), which is also recalled in the source paper.

Before this upgrade, the run indexes were searched for arXiv:2504.03386, the
source title, and core phrases around order-preserving isometric embeddings,
unique positive isometry extensions, irreducible surjections, Banach-Stone, and
Holsztynski.  The only exact target hit was the prior partial packet
\texttt{2504.03386\_peakpoint\_unique\_positive\_isometry\_extension}.  A
bounded web search for exact phrases from the source problem and close variants
did not find a direct later answer.  The theorem should therefore be reviewed
as a new literal-answer packet, with the caveat that it gives conditions rather
than a full classification.

\section*{Human-review recommendation}
Review as a candidate full literal answer to the source problem.  The main
points to check are: (i) the source context justifies the unital hypothesis,
(ii) Holsztynski's theorem is being applied in the correct real \(C(K)\) into
isometry form, and (iii) the phrase "find conditions" is accepted as answered
by the irreducible-surjection condition rather than requiring a classification
of all possible embeddings.

\begin{thebibliography}{9}
\bibitem{PaulRao}
Tanmoy Paul and T. S. S. R. K. Rao,
\emph{Order-preserving unique Hahn-Banach extensions},
arXiv:2504.03386, 2025.
\url{https://arxiv.org/abs/2504.03386}

\bibitem{Holsztynski}
W. Holsztynski,
\emph{Continuous mappings induced by isometries of spaces of continuous functions},
Studia Mathematica 26 (1966), 133--136.
\end{thebibliography}

\end{document}
